\documentclass[a4paper, 12pt]{article}

\usepackage[brazil]{babel}
\usepackage[utf8]{inputenc}
\usepackage{amsmath}
\usepackage{indentfirst}
\usepackage{graphicx}
\usepackage{multicol}
\usepackage{url}
\usepackage{hyperref}
\usepackage{booktabs}

\hypersetup{
    colorlinks=true,
    linkcolor=blue,
    citecolor=blue,
    urlcolor=blue
}

\begin{document}

% % % % % % % % % % % % % % % CAPA % % % % % % % % % % % % % % %
\begin{titlepage}
	\begin{center}
		\textbf{\large UNIVERSIDADE DE BRASÍLIA -- UnB}\\
		\textbf{\large FACULDADE DE CIÊNCIAS E TECNOLOGIAS EM ENGENHARIA -- FCTE}\\
		\textbf{\large CURSO DE ENGENHARIA DE SOFTWARE}\\
		\large DISCIPLINA: QUALIDADE DE SOFTWARE 1 -- TURMA 02\\
		\large PROFESSORA: Ms. CRISTIANE SOARES RAMOS\\
		\vspace{40pt}
		
		\textbf{\large EQUIPE KAREN JONES}
		
		\vspace{80pt}
		
		\textbf{\LARGE{AVALIAÇÃO DA QUALIDADE DE PRODUTO DE SOFTWARE}}\\
		\vspace{12pt}
		\textbf{\Large{Estudo de Caso no Software Open Source Excalidraw segundo a ISO/IEC 25000 (SQuaRE)}}\\
		
		\vspace{\fill}
		
		\begin{flushleft}
			\textbf{Integrantes da Equipe:}
			\begin{itemize}
				\itemsep 2pt
				\item Christopher
				\item Davi Pierre
				\item Gabriel Monteiro
				\item Letícia de Carvalho
				\item Luiz Henrique Soares
				\item Maria Eduarda
			\end{itemize}
		\end{flushleft}
		
		\vspace{30pt}
		
		Gama -- DF\\
		2026
	\end{center}
\end{titlepage}

% % % % % % % % % % % % % FOLHA DE ROSTO % % % % % % % % % % % %
\begin{titlepage}
	\begin{center}
		\textbf{\large UNIVERSIDADE DE BRASÍLIA -- UnB}\\
		\textbf{\large FACULDADE DE CIÊNCIAS E TECNOLOGIAS EM ENGENHARIA -- FCTE}\\
		\textbf{\large ENGENHARIA DE SOFTWARE}\\
		\vspace{40pt}
		
		\textbf{\large EQUIPE KAREN JONES}
		
		\vspace{60pt}
		
		\textbf{\LARGE{AVALIAÇÃO DA QUALIDADE DE PRODUTO DE SOFTWARE}}\\
		\vspace{12pt}
		\textbf{\Large{Estudo de Caso no Software Open Source Excalidraw segundo a ISO/IEC 25000 (SQuaRE)}}\\
	\end{center}
	
	\vspace{30pt}
	
	\begin{flushright}
		\begin{minipage}{0.58\textwidth}
			\small
			Trabalho técnico e artigo apresentado como requisito parcial de avaliação na disciplina de \textbf{Qualidade de Software 1} (FGA0315), no 2º semestre de 2026, ministrada pela Profa. Dra. Cristiane Soares Ramos, no curso de Engenharia de Software da Faculdade de Ciências e Tecnologias em Engenharia (FCTE) da Universidade de Brasília (UnB).
			\vspace{10pt}
			
			\textbf{Requisitante:} Equipe Karen Jones (QA Independente)\\
			\textbf{Produto Avaliado:} Excalidraw\\
			\textbf{Repositório do Projeto:} \url{https://github.com/FCTE-Qualidade-de-Software-1/2026-2_T02_KAREN_JONES}
		\end{minipage}
	\end{flushright}
	
	\vspace{\fill}
	
	\begin{center}
		Gama -- DF\\
		2026
	\end{center}
\end{titlepage}

\newpage
\tableofcontents
\thispagestyle{empty}

\newpage
\pagenumbering{arabic}

% % % % % % % % % % % % SEÇÕES DO ARTIGO % % % % % % % % % % %
\section{Resumo}
Este trabalho apresenta a avaliação da qualidade de produto do software livre Excalidraw, fundamentada nas normas da família ISO/IEC 25000 (SQuaRE). A avaliação foi conduzida sob o papel de equipe de garantia da qualidade (QA) independente, aplicando o processo formal da ISO/IEC 25040 e selecionando três características do modelo ISO/IEC 25010:2023: Manutenabilidade, Confiabilidade e Flexibilidade. As métricas foram desdobradas através da abordagem Goal-Question-Metric (GQM), e a gestão do projeto foi guiada por indicadores do modelo PSM/CID.

\vspace{10pt}
\noindent\textbf{Palavras-chave:} Qualidade de Software, SQuaRE, ISO/IEC 25010, ISO/IEC 25040, GQM, Excalidraw.

\newpage
\section{Introdução}
O Excalidraw é uma ferramenta web de código aberto amplamente utilizada para diagramação e colaboração em tempo real. Desenvolvido em TypeScript e React com arquitetura monorepo, o projeto conta com expressiva adoção comunitária e acadêmica.

O objetivo desta avaliação é \textbf{indicar pontos para melhoria no produto}, apoiando desenvolvedores, mantenedores e integradores da biblioteca \texttt{@excalidraw/excalidraw}. A equipe Karen Jones atua como requisitante independente de QA, sem vínculo formal prévio com os mantenedores do software.

\subsection{Link de Rastreabilidade do Repositório}
Todos os artefatos de medição, atas de reunião, relatórios PSM/CID e códigos-fonte da documentação estão disponíveis no repositório oficial da equipe:
\begin{center}
    \url{https://github.com/FCTE-Qualidade-de-Software-1/2026-2_T02_KAREN_JONES}
\end{center}

\section{Referencial Teórico}
Esta seção aborda as normas e métodos utilizados:
\begin{itemize}
    \item \textbf{Família ISO/IEC 25000 (SQuaRE):} Conjunto de normas que padroniza os requisitos e a avaliação de qualidade de produtos de software.
    \item \textbf{ISO/IEC 25040 (Processo de Avaliação):} Roteiro estruturado em quatro fases (Requisitos, Especificação, Projeto e Execução).
    \item \textbf{ISO/IEC 25010:2023 (Modelo de Qualidade do Produto):} Versão atualizada da norma, adotando características como Manutenabilidade, Confiabilidade e Flexibilidade.
    \item \textbf{Goal-Question-Metric (GQM):} Metodologia de definição orientada a objetivos para métricas de software.
    \item \textbf{PSM/CID (Practical Software Measurement):} Acompanhamento gerencial do processo, cronograma e esforço.
\end{itemize}

\section{Fase 1: Requisitos de Avaliação}
Na primeira fase do processo de avaliação da ISO/IEC 25040, foram estabelecidos o escopo, o propósito da avaliação e o alinhamento com as partes interessadas.

\subsection{Propósito e Requisitante da Avaliação}
A equipe Karen Jones declarou-se como requisitante no papel de equipe de QA independente. O uso pretendido dos resultados compreende o fornecimento de um diagnóstico técnico transparente e a elaboração de recomendações e propostas de melhoria (incluindo prototipação de alta fidelidade para o site).

\subsection{Partes Interessadas e Critérios de Sucesso}
Foram mapeados os seguintes grupos de interessados: usuários finais, desenvolvedores do projeto, mantenedores, contribuidores open source, integradores da biblioteca, equipes de implantação/operação, equipes de qualidade, equipe avaliadora/docentes e comunidade adotante. Para cada grupo, foram delimitados critérios de sucesso específicos e sua influência nas características de qualidade selecionadas.

\subsection{Descrição Estruturada do Produto e Arquitetura}
O Excalidraw é classificado tecnicamente como uma aplicação web rich client (PWA/SPA) associada a uma biblioteca gráfica reutilizável, estruturada em monorepo com TypeScript e React. O repositório divide-se entre a aplicação de produção (\texttt{excalidraw-app/}), o pacote de biblioteca distribuído via npm (\texttt{packages/excalidraw/}) e módulos utilitários de suporte geométrico (\texttt{packages/math/} e \texttt{packages/utils/}). As interfaces externas compreendem o motor de renderização Rough.js (traçado estilo esboço), as APIs 2D Canvas do navegador, mecanismos de persistência local (IndexedDB e LocalStorage) e o servidor de colaboração síncrona (\textit{Room Server}) com criptografia ponta a ponta. 

Em termos metodológicos para a avaliação, a arquitetura delimita claramente o que é viável medir: a qualidade interna, acoplamento e cobertura de testes são integralmente auditáveis via análise estática, enquanto subcaracterísticas dependentes de execução (Disponibilidade, Tolerância a falhas e Recuperabilidade) são investigadas através de inspeção estática das rotinas de persistência e tratamento de exceções, contornando a indisponibilidade de telemetria fechada de produção.

\subsection{Seleção das Características e Método de Priorização}
A seleção das características não foi arbitrária: aplicou-se um método formal ponderando três critérios — abrangência de partes interessadas (quantos grupos elegem a característica como prioritária), viabilidade técnica de medição e impacto nas metas de melhoria do software. Foram selecionadas três características da ISO/IEC 25010:2023:
\begin{enumerate}
    \item \textbf{Manutenabilidade (Prioridade 1):} Prioritária para 4 dos 7 grupos de interessados, avaliada em Modularidade, Reusabilidade, Analisabilidade, Modificabilidade e Testabilidade.
    \item \textbf{Confiabilidade (Prioridade 2):} Prioritária para 3 grupos, avaliada em Maturidade, Disponibilidade, Tolerância a falhas e Recuperabilidade.
    \item \textbf{Flexibilidade (Prioridade 3):} Prioritária para 2 grupos, avaliada em Adaptabilidade, Escalabilidade, Instalabilidade e Substituibilidade.
\end{enumerate}

A seleção envolveu \textit{trade-offs} deliberados: priorizou-se a profundidade analítica em detrimento da cobertura superficial das nove características da norma (excluindo Usabilidade, Desempenho e Segurança), além de balancear o esforço de coleta estática para viabilizar conclusões rastreáveis no semestre letivo.

\subsection{Relação com os Objetivos de Desenvolvimento Sustentável (ODS)}
O Excalidraw foi relacionado aos seguintes ODS da Agenda 2030 da ONU:
\begin{itemize}
    \item \textbf{ODS 4 (Educação de Qualidade):} Ferramenta gratuita de apoio pedagógico e colaboração (Metas 4.3, 4.4 e 4.a).
    \item \textbf{ODS 9 (Indústria, Inovação e Infraestrutura):} Fomento à tecnologia aberta e suporte a funcionamento offline (Metas 9.5, 9.b e 9.c).
    \item \textbf{ODS 10 (Redução das Desigualdades):} Redução de barreiras financeiras e adoção de formatos abertos de exportação (Metas 10.2 e 10.3).
\end{itemize}

\section{Fase 2: Especificação da Avaliação}
Aplicação da abordagem Goal-Question-Metric (GQM) para derivar questões e métricas a partir dos objetivos de qualidade, acompanhada dos níveis de pontuação e critérios de julgamento.

\section{Fase 3: Planejamento da Avaliação}
Definição dos métodos de medição, ferramentas de análise estática e dinâmica, especificação do ambiente de testes e cronograma de execução da avaliação.

\section{Fase 4: Execução da Avaliação e Ações de Melhoria}
Apresentação dos dados coletados, respostas ao GQM, julgamento dos resultados segundo os critérios pré-estabelecidos e elaboração de ações de melhoria (com desenvolvimento de protótipo de alta fidelidade).

\section{Gestão do Projeto (Métricas PSM/CID)}
Monitoramento contínuo da equipe através de métricas de Desempenho dos Processos, Calendário e Progresso, e Recursos e Custos, registradas nas atas e relatórios de medição.

\section{Considerações Finais}
Síntese das lições aprendidas, limitações do estudo e impacto da avaliação para a sustentabilidade do Excalidraw.

\newpage
\addcontentsline{toc}{section}{Referências}
\begin{thebibliography}{99}

\bibitem{iso25010}
ISO/IEC. \textit{ISO/IEC 25010:2023 -- Systems and software engineering -- Systems and software Quality Requirements and Evaluation (SQuaRE) -- Product quality model}. International Organization for Standardization, Geneva, 2023.

\bibitem{iso25040}
ISO/IEC. \textit{ISO/IEC 25040:2011 -- Systems and software engineering -- Systems and software Quality Requirements and Evaluation (SQuaRE) -- Evaluation process}. International Organization for Standardization, Geneva, 2011.

\bibitem{basili1994gqm}
BASILI, V. R.; CALDIERA, G.; ROMBACH, H. D. \textit{The Goal Question Metric Approach}. Encyclopedia of Software Engineering, John Wiley \& Sons, 1994.

\bibitem{psm}
CARD, D. et al. \textit{Practical Software Measurement: Objective Information for Decision Makers}. Department of Defense, 2000.

\bibitem{onu2015ods}
ONU. \textit{Transformando Nosso Mundo: A Agenda 2030 para o Desenvolvimento Sustentável}. Organização das Nações Unidas, 2015.

\end{thebibliography}

\newpage
\appendix
\section{Tabela de Contribuição dos Membros da Equipe}
Em conformidade com as premissas da disciplina de Qualidade de Software 1, a Tabela~\ref{tab:contribuicao} apresenta o esforço e a participação percentual individual dos integrantes para a entrega do projeto.

\begin{table}[!ht]
\centering
\caption{Contribuição individual dos membros da Equipe Karen Jones}
\label{tab:contribuicao}
\vspace{6pt}
\begin{tabular}{lcccc}
\toprule
\textbf{Integrante} & \textbf{Issues} & \textbf{Horas Estimadas} & \textbf{Horas Reais} & \textbf{Esforço (\%)} \\
\midrule
Christopher          & 5 & 0.8h & 2.5h & 38.46\% \\
Davi Pierre          & 6 & 0.0h & 0.0h & 0.00\% \\
Gabriel Monteiro     & 6 & 0.0h & 0.0h & 0.00\% \\
Letícia de Carvalho  & 6 & 0.0h & 0.0h & 0.00\% \\
Luiz Henrique Soares & 6 & 0.0h & 0.0h & 0.00\% \\
Maria Eduarda        & 6 & 2.8h & 4.0h & 61.54\% \\
\midrule
\textbf{Total}       & \textbf{35} & \textbf{3.6h} & \textbf{6.5h} & \textbf{100.00\%} \\
\bottomrule
\end{tabular}
\end{table}

\section{Declaração de Uso de Inteligência Artificial}
Em atendimento aos critérios de avaliação da disciplina, declara-se que ferramentas de Inteligência Artificial generativa foram utilizadas para auxiliar nas seguintes atividades:
\begin{itemize}
    \item \textbf{Ferramenta utilizada:} Antigravity IDE (Claude / Gemini).
    \item \textbf{Finalidade:} Auxílio na estruturação de templates LaTeX, formatação de tabelas, revisão gramatical e levantamento de metas e indicadores dos ODS da Agenda 2030 da ONU.
    \item \textbf{Supervisão e Validação Humana:} Todo o conteúdo textual gerado ou revisado por ferramentas automatizadas foi minuciosamente conferido, ajustado e validado pelos membros da equipe, garantindo a conformidade técnica, ausência de jargões genéricos ou inconsistências conceituais.
\end{itemize}

\end{document}
